% Содержимое отчета по курсу Анализ алгоритмов

\aaunnumberedsection{ВВЕДЕНИЕ}{sec:intro}

\textbf{Параллелизм} описывает одновременное выполнение различных последовательностей действий~\cite{posix-threads}. Следовательно, параллельные вычисления подразумевают выполнение операций одновременно. Распараллеливание вычислений может повысить временную эффективность программы на многопроцессорных системах и даже на однопроцессорных, если программа часто простаивает, ожидая ввода-вывода~\cite{tanenbaum}.

Некоторые программы можно разделить на отдельные компоненты, каждая из которых выполняет определённую операцию над данными и передаёт результат следующему компоненту. В результате формируется конвейер обработки данных, где элементы этого конвейера могут функционировать параллельно друг с другом.

Цель данной работы -- разработка программного обеспечения для обработки файлов страниц с сайта vkusnoprigotovim.ru.

Задачи работы включают разработку ПО, выполняющего обработку файлов рецептов в конвейере и исследование временных характеристик разработанной программы.

\clearpage
\aasection{Входные и выходные данные}{sec:input-output}

Входными данными для программы является набор файлов с информацией о рецептах.
Выходными данными является информация о рецептах, записанная в хранилище.

\clearpage
\aasection{Преобразование входных данных в выходные}{sec:algorithm}

Для реализации ПО был использован язык Golang~\cite{go} с использованием корутин для организации конвейра.

Конвейер состоит из 3-х частей:

\begin{enumerate}
	\item чтение данных из файла;
	\item извлечение необходимого подмножества данных;
	\item запись извлеченных данных в хранилище.
\end{enumerate}

Каждая из частей запускается в отдельной корутине~\cite{go-mem}, которые передают данные через каналы~\cite{go-mem}.

Элементом данных выступает структура, которая хранит в себе информацию о рецепте, информацию о моментах времени начала и конца обработки отдельной частью программы и метаинформацию для обработчиков конвейера. Структура представлена в листинге~\ref{unit}.

\begin{lstlisting}[label=unit,caption={Данные, передающиеся между частями конвейера},language=go]
type Task struct {
	ID              int
	FileName        string
	Text            string
	Data            Data
	CreationTime    time.Time
	QueueTimes      []time.Time
	ProcessingTimes []time.Time
	CompletionTime  time.Time
	Context         *TaskContext
}

type Data struct {
	Id          int          `db:"id"`
	Url         string       `db:"url"`
	Title       string       `db:"title"`
	Ingredients []Ingredient `db:"ingredients"`
	Steps       []string     `db:"steps"`
	ImageUrl    string       `db:"image_url"`
}

type Ingredient struct {
	Name     string `json:"name"`
	Unit     string `json:"unit"`
	Quantity string `json:"quantity"`
}

type TaskContext struct {
	Dsn string
}
\end{lstlisting}

Полученные в результате работы программы рецепты сохраняются в базе данных PostgreSQL~\cite{pgs}, а из неё отдельным скриптом моугт быть сохранены в виде json-файлов.

\clearpage
\aasection{Примеры работы программы}{sec:demo}

В листинге~\ref{lst:exmpl} представлен пример рецепта, обработанного программой.

\begin{lstlisting}[label=lst:exmpl,caption={Данные рецепта}]
{
"id": 498,
"issue_id": 9170,
"url": "https://vkusnoprigotovim.ru/zagotovki/zharenye-baklazhany",
"title": "Жареные баклажаны на зиму",
"ingredients": [
    {
        "name": "баклажаны - по вкусу",
        "unit": "",
        "quantity": ""
    }
],
"steps": [
    "Баклажаны почистить и порезать длинненькими кусочками.",
    "Через 2 часа воду слить и промыть баклажаны.",
    "Обжаренные баклажаны выложить на салфетку и положить в морозильник.",
    "Зимой это хороший способ разнообразить свое меню."
],
"image_url": "/zagotovki/zharenye-baklazhany-na-zimu/zharenye-baklazhany.jpg"
}
\end{lstlisting}

\clearpage
\aasection{Тестирование}{sec:tests}

Тестирование программы проводилось загрузкой в качестве входных данных 500 файлов различных страниц рецептов. При этом отслеживалось количество удачных обработок данных, а также выборочная ручная проверка 10 случайных рецептов. 

Тестирование было успешно пройдено.

\clearpage
\aasection{Описание исследования}{sec:study}

Технические характеристики устройства,на котором проводились замеры:

\begin{itemize}
	\item процессор: AMD Ryzen 7 5800U (16) @ 4.51 ГГц;
	\item оперативная память: 16 Гб;
	\item ядро ОС Linux 6.11.6-zen1-1-zen.
\end{itemize}

При проведении замеров ноутбук был включён в сеть и были запущены только системные приложения.

На вход программе были поданы 500 файлов, каждый из которых был успешно обработан.

По итогу выполнения программы выводятся:

\begin{itemize}
    \item среднее время существования задачи;
    \item среднее время ожидания задачи в каждой из очередей;
    \item среднее время обработки задачи на каждой из стадий.
\end{itemize}

В листинге~\ref{log} представлены результаты замеров.

\begin{lstlisting}[label=log,caption={Данные журналирования}]
Average task lifetime: 1.379249857s
Average queue time for Stage 1: 323.860945ms
Average queue time for Stage 2: 458.44184ms
Average queue time for Stage 3: 590.42342ms
Average processing time for Stage 1: 31.578mks
Average processing time for Stage 2: 66.874mks
Average processing time for Stage 3: 6.490491ms
\end{lstlisting}

Самым долгим этапом обработки файлов оказался этап преобразования данных. Самым быстрым - этап записи в хранилище.

\clearpage
\aaunnumberedsection{ЗАКЛЮЧЕНИЕ}{sec:outro}

Было разработано программное обеспечение, выполняющее обработку файлов страниц с сайта vkusnoprigotovim.ru.

В ходе работы были решены следующие задачи:

\begin{itemize}
	\item разработка ПО, выполняющего обработку файлов рецептов в конвейере;
    \item исследование временных характеристик разработанной программы.
\end{itemize}

В результате исследования ПО, наиболее продолжительным этапом обработки файлов оказался этап преобразования данных, а самым быстрым - этап записи в базу данных.
